\section{Infrastructure Design}
\label{sec:system}
\subsection{Bottleneck Analysis}

The main challenge in \textbf{SSO} implementation comes from the \textbf{bracket-and-bisect} root solver that runs at every update. To satisfy the tangent-space constraint, we must find a Lagrange multiplier $\lambda$ such that $h(\lambda)=0$. We first expand the search interval by \textbf{bracketing}, and then we run \textbf{bisection} inside the bracket to obtain a $\lambda^{\star}$ that meets a target tolerance. This solver introduces non-trivial overhead:
\begin{enumerate}[left=0pt]
    \item \textbf{Workload Imbalance:} different bracketing range and tolerance settings can change the number of search and bisection steps, leading to unstable runtime and workload imbalance between devices.
    \item \textbf{Computational Cost:} Each step in the search also evaluates $h(\lambda)$, which calls $\operatorname{msign}(\widehat{\mM}+\lambda\mTheta)$; these extra matrix computations accumulate and add noticeable cost to every optimizer step.
    \item \textbf{Synchronization Overhead:} Iterative search creates frequent synchronization between the GPU and the CPU, because the algorithm must finish each evaluation and then check a scalar condition before choosing the next $\lambda$ and continuing.
\end{enumerate}

\subsection{Optimization Pipeline}

We introduce a holistic optimization pipeline designed to mitigate these overheads while preserving numerical precision. \textbf{All performance statistics are collected from Dense 1.7B model pretraining.}

\paragraph{Atomic Module Sharding.}
We employ a fine-grained, parameter-wise sharding strategy \citep{emergeing_optimizer}. As noted by \citep{liu2025muonscalablellmtraining}, while standard ZeRO-1 is efficient for element-wise optimizers (e.g. AdamW), its flat-buffer sharding approach is incompatible with spectral methods that require full gradient matrices to compute updates. To reconcile this, we shard parameters as \emph{atomic modules} rather than flattened buffers. An atomic module is defined as the minimal independent weight matrix required to remain intact for spectral operations (see~\cref{sec:module_gran}).

\paragraph{Load Balancing Strategy.}
To resolve the workload imbalance caused by variable solver depths, we employ a ``ping-pong'' load balancing strategy adapted from \citep{emergeing_optimizer}. Empirical results indicate this method outperforms both greedy size-descent sorting and default round-robin allocation. We sort atomic modules by size and assign them to DP ranks in an alternating zigzag pattern. This heuristic effectively balancing the solver workload without complex runtime scheduling.
Finally, we synchronize the updated params using iterative All-Gather collective from \cite{emergeing_optimizer}.
\paragraph{Adaptive Kernel Selection.}
We observe that the optimal implementation for Matrix Sign computation is highly sensitive to matrix dimensions. As shown in~\cref{tab:perf_breakdown}, applying specialized kernels indiscriminately can degrade performance. We therefore implement an Adaptive Dispatcher:
\begin{itemize}[left=0pt]
    \item \textbf{Small Matrices (\(<512\))}: We use a JIT-compiled PyTorch implementation built on \texttt{torch.addmm}. This avoids the launch overhead of specialized kernels.
    \item \textbf{Large Matrices ($\ge512$)}: We dispatch to a custom Triton kernel implementing the SYmmetric Rank-K (SYRK)  \cite{emergeing_optimizer} update, which exploits the symmetry of Newton--Schulz iterations to halve memory reads.
\end{itemize}

\paragraph{Multi-Stream Parallelism.}
For layers composed of many small independent matrices (e.g. per head attention split), single-stream execution suffers from kernel launch latency bubbles. We exploit this independence by dispatching spectral updates across multiple CUDA streams.

\paragraph{Mixed-Precision.}
The Power Iteration for spectral norm estimation is performed in BFloat16, while the sensitive \textbf{msign remains in FP32 with 8 iterations}.

\paragraph{Cache Top Singular Vectors.}
The singular vectors of model weights evolve slowly during training. Leveraging this temporal locality, we initialize the current Power Iteration using the cached singular vectors $u$ and $v$ from the previous step. This reuse mechanism drastically accelerates convergence, requiring only a few iterations to maintain approximation accuracy.

\begin{table}[hbt!]
\centering
\caption{Optimization breakdown on end-to-end latency
for 4M tokens/step on NVIDIA B200.
\goodarrow~denotes improvement, while~\badarrow~denotes regression. Note there is still room for improvement.}
\label{tab:perf_breakdown}
\resizebox{0.9\columnwidth}{!}{

\begin{tabular}{l r r r}
\toprule
   & \multicolumn{1}{c}{Time (ms)} & \multicolumn{1}{c}{$\Delta$ vs Baseline} & \multicolumn{1}{c}{$\Delta$ vs Prev.} \\
% Fully Optimized & 7666.32  & -3262.16 (-29.85\%)~\goodarrow & -     \\
\midrule
Naive Baseline (No opt.) & 10928.5  & \multicolumn{1}{c}{-}                              & \multicolumn{1}{c}{-}                           \\
+ Load balance \& All Gather   & 9365.5 & -1563.0 (-14.3\%)~\goodarrow & -1563.0 (-14.3\%)~\goodarrow \\
+ Triton SYRK Kernel      & 10284.2  & -644.3 \hspace{0.5em}(-5.9\%)~\goodarrow   & +918.7 \hspace{0.3em}(+9.8\%)~\badarrow \\
+ Adaptive \& Multi-stream       & 9383.4   & -1545.1 (-14.2\%)~\goodarrow & -900.8 \hspace{0.5em}(-8.8\%)~\goodarrow \\
+ BF16 \& Torch.compile (\textbf{Final}) & \textbf{7666.3}   & \textbf{-3262.2 (-29.9\%)}~\goodarrow & \textbf{-1717.1 (-18.3\%)}~\goodarrow \\
\bottomrule
\end{tabular}%
}
\end{table}

\begin{table}[ht]
\centering
\caption{End-to-end per-step latency for 4M tokens/step on NVIDIA B200.  \textcolor{MuonColor}{Muon} as baseline.}
\label{tab:perf_end2end}
\resizebox{0.9\columnwidth}{!}{%
\begin{tabular}{ccccc}
\toprule
 & \textcolor{AdamWColor}{AdamW}
 & \textcolor{MuonColor}{Muon}
 & \textcolor{MuonSphereColor}{MuonSphere}
 & \textcolor{SpectralSphereColor}{Spectral Sphere} \\
\midrule
Time (ms) & 6734.15 (-2.10\%) & 6878.83 & 6949.85 (+1.03\%) & 7666.32 (+11.45\%) \\
\bottomrule
\end{tabular}%
}
\end{table}

\subsection{Future Improvements}
\label{sec:infra_future}
\begin{itemize}[left=0pt]
    \item \textbf{GPU-Native Solver.}
    Profiling indicates that the current CPU-based bisection solver introduces latency due to frequent device-host synchronization. Although the bracketing phase converges rapidly ($<2$ steps), the bisection phase averages 5--7 steps, creating synchronization bubbles. Future work will prioritize implementing a pure GPU-native solver to eliminate these overheads. Additionally, we plan to adopt higher convergence algorithms,  such as Brent's method or $n$-section search, and optimize initial bracketing intervals to minimize msign calls.

    \item \textbf{Kernel Optimization.}
    The theoretical advantage of SSO relies on the accuracy of the tangent space projection; when errors are high, the update direction may degrade to the ``worst-case'' Muon update. Currently, we ensure precision via 8 iterations of msign in FP32. To follow standard Muon practices, we plan to use msign in BFloat16 within 5 steps. Furthermore, we intend to replace current JIT-compiled operations with custom, fully optimized kernels (e.g. for batched msign and Power Iterations) to better exploit the hardware features of next-generation GPUs.

    \item \textbf{Low-precision Training.}
    While weight matrics are spectral constrained, we find the residual stream remains the primary source of outliers. We aim to explore ``fully manifold constrained'' architectures, i.e. using mHC \citep{xie2025mhcmanifoldconstrainedhyperconnections}. Additionally, given SSO's demonstrated stability, we plan to explore low-precision training (e.g. FP8/NVFP4) to leverage the high throughput of low-bit arithmetic for better training efficiency.
\end{itemize}
